\chapter{Repositório do projeto}
\label{ap:repositorio}

Este apêndice descreve a organização do repositório e os comandos necessários
para reproduzir a execução, os testes e as evidências.

\section*{Organização dos diretórios}
\label{ap:arvore}

\begin{verbatim}
lab404/
|-- assets/            # GLB exportado do Blender e lista de assets
|-- artigo/            # artigo cientifico do projeto (LaTeX)
|-- audio/             # efeitos sonoros originais (gerados por script)
|-- config/            # mapa de entrada e personalidade da IA
|-- design/            # conceito visual e level design
|-- docs/              # documentacao, imagens e videos de evidencia
|-- monografia/        # esta monografia (LaTeX + PDF)
|-- scenes/            # cenas do jogo (.tscn)
|-- scripts/           # GDScript do jogo + adaptador de IA (Python)
|-- tests/             # suite automatizada (motor + Python)
|-- tools/             # gerador de audio
|-- start.sh / stop.sh # ligar e desligar jogo e servico
|-- SPECIFICATION.md   # requisitos FR/NFR e criterios de aceitacao
`-- STATUS.md / TASKS.md / HANDOFF.md / AGENTS.md
\end{verbatim}

\section*{Comandos de reprodução}
\label{ap:comandos}

\begin{itemize}
  \item \textbf{Verificação funcional completa:} \texttt{tests/run\_all.sh}
  (executa as cenas de teste do motor e os testes do adaptador);
  \item \textbf{execução do jogo:} \texttt{godot4 --path .} ou
  \texttt{./start.sh} (que também inicia o serviço de IA, quando o ambiente
  Python existe);
  \item \textbf{serviço de IA:} \texttt{.venv/bin/uvicorn
  scripts.aria\_adapter:app --port 8000}, com o modelo local servido pelo
  Ollama;
  \item \textbf{inspeção do asset 3D:} script de contagem de nós, materiais e
  colisões do arquivo GLB;
  \item \textbf{evidências visuais:} execução da cena de captura com as
  variáveis de pose e de conclusão da missão;
  \item \textbf{compilação desta monografia:}
  \texttt{monografia/build.sh} (pdflatex, bibtex e duas passagens
  adicionais).
\end{itemize}

\section*{Documentos de controle do projeto}
\label{ap:controle}

O desenvolvimento manteve cinco documentos de estado que registram, de forma
contínua, o que está pronto, o que está pendente e quais são as evidências:
especificação síncrona (\texttt{SPECIFICATION.md}), estado
(\texttt{STATUS.md}), tarefas (\texttt{TASKS.md}), transferência entre agentes
(\texttt{HANDOFF.md}) e regras de projeto (\texttt{AGENTS.md}). A monografia
mantém, no diretório \texttt{monografia/}, cinco arquivos de controle que
registram o estado desta produção acadêmica:

\begin{itemize}
  \item \texttt{MONOGRAFIA\_STATUS.md};
  \item \texttt{MONOGRAFIA\_PLANO.md};
  \item \texttt{MONOGRAFIA\_EVIDENCIAS.md};
  \item \texttt{MONOGRAFIA\_RASTREABILIDADE.md};
  \item \texttt{MONOGRAFIA\_PENDENCIAS.md}.
\end{itemize}

\section*{Disponibilidade}
\label{ap:disponibilidade}

O código, a especificação, as capturas e o vídeo de demonstração estão
publicamente disponíveis \cite{lab404repo,lab404video}. As imagens usadas
nesta monografia foram copiadas do diretório de evidências do projeto; os
diagramas deste documento possuem código-fonte versionado (formato Mermaid
\cite{mermaid2026}) e são regeneráveis por comando.
